%! Absolute Value Functions \& Transformations — Unit 2, Lesson 2.3 — Group Activity (Answer Key)
\documentclass[10pt]{article}
\usepackage{algebra2-article}
\usepackage{algebra2-key}   % pulls in -boxes; do NOT also load -boxes
\newcommand{\TallMath}[1]{$\displaystyle #1\rule[-1.4em]{0pt}{3.2em}$}
% Coordinate grid: {xmin}{xmax}{ymin}{ymax}
\newcommand{\lgrid}[4]{%
  \draw[step=1, linegray!40, very thin] (#1,#3) grid (#2,#4);
  \draw[->, semithick, charcoal] (#1,0)--(#2,0) node[below right, font=\scriptsize]{$x$};
  \draw[->, semithick, charcoal] (0,#3)--(0,#4) node[left, font=\scriptsize]{$y$};}

\begin{document}

\pageheader{Unit 2, Lesson 2.3}{Group Activity --- Answer Key}
\namepartnerperiod

\vspace{0.08in}

\begin{headlinebox}{forestbg}
\small\textbf{Reading \& Building V's.} Every absolute value graph is the parent $y=|x|$ moved and
stretched. Work with your partner. In \textbf{vertex form} $g(x)=a|x-h|+k$, the vertex is $(h,k)$, the
axis is $x=h$, and $a$ tells you \emph{up/down} (sign) and \emph{narrow/wide} ($|a|$). Tip: the vertex
sits where the inside of the bars equals zero.
\end{headlinebox}

\vspace{0.06in}

\begin{tcolorbox}[colback=white, colframe=black!40, title=\textbf{Tier R --- Remediate},
    fonttitle=\bfseries, arc=1.5mm, left=3mm, right=3mm, top=2mm, bottom=2mm, breakable]
\begin{enumerate}[leftmargin=1.7em, itemsep=9pt]
  \item \textbf{Complete the table} for $y=|x|+1$, then give the vertex.
        \begin{center}
        \renewcommand{\arraystretch}{1.2}
        \begin{tabular}{|c|c|c|c|c|c|}
        \hline
        $x$ & $-2$ & $-1$ & $0$ & $1$ & $2$ \\ \hline
        $y$ & ${\color{keyred}\mathbf{3}}$ & ${\color{keyred}\mathbf{2}}$ & ${\color{keyred}\mathbf{1}}$ & ${\color{keyred}\mathbf{2}}$ & ${\color{keyred}\mathbf{3}}$ \\ \hline
        \end{tabular}
        \end{center}
        Vertex: $({\color{keyred}\mathbf{0}},\,{\color{keyred}\mathbf{1}})$.
  \item \textbf{Read the graph.} The V below is one absolute value function.
        \begin{center}
        \begin{tikzpicture}[scale=0.5]
          \lgrid{-1}{5}{-1}{5}
          \draw[navy, dashed, semithick] (2,-1)--(2,5);
          \draw[very thick, forest] (-1,4)--(2,1)--(5,4);
          \fill[charcoal] (2,1) circle (3.5pt);
        \end{tikzpicture}
        \end{center}
        Vertex $({\color{keyred}\mathbf{2}},\,{\color{keyred}\mathbf{1}})$; \; axis $x={\color{keyred}\mathbf{2}}$;
        \; minimum value \ans{$1$}; \; range \ans{$y\ge 1$}.
  \item \textbf{Match} each equation to its graph (write A or B).
        \begin{center}
        \begin{tikzpicture}[scale=0.42]
          \begin{scope}
            \lgrid{-1}{5}{-1}{4}
            \draw[very thick, forest] (-1,3)--(2,0)--(5,3);
            \node[charcoal, font=\scriptsize] at (2,-2.2){\textbf{A}};
          \end{scope}
          \begin{scope}[xshift=9cm]
            \lgrid{-3}{3}{-1}{6}
            \draw[very thick, forest] (-3,5)--(0,2)--(3,5);
            \node[charcoal, font=\scriptsize] at (0,-2.2){\textbf{B}};
          \end{scope}
        \end{tikzpicture}
        \end{center}
        $y=|x-2|$ is graph \ans{A}. \qquad $y=|x|+2$ is graph \ans{B}.
\end{enumerate}
\end{tcolorbox}

\begin{tcolorbox}[colback=white, colframe=black!40, title=\textbf{Tier A --- Approaching Proficiency},
    fonttitle=\bfseries, arc=1.5mm, left=3mm, right=3mm, top=2mm, bottom=2mm, breakable]
\begin{enumerate}[leftmargin=1.7em, itemsep=10pt]
  \item \textbf{Name the features} of each function.
        \begin{center}
        \renewcommand{\arraystretch}{1.4}
        \begin{tabularx}{\linewidth}{@{}l X X X X@{}}
        & \textbf{Vertex} & \textbf{Axis} & \textbf{Opens} & \textbf{Max/Min value} \\
        (a) $y=|x-2|+1$ & \ans{$(2,1)$} & \ans{$x=2$} & \ans{up} & \ans{min $1$} \\
        (b) $y=-|x+3|$   & \ans{$(-3,0)$} & \ans{$x=-3$} & \ans{down} & \ans{max $0$} \\
        (c) $y=2|x|-4$   & \ans{$(0,-4)$} & \ans{$x=0$} & \ans{up} & \ans{min $-4$} \\
        \end{tabularx}
        \end{center}
  \item \textbf{Match} each equation to a graph (write I, II, or III).
        \begin{center}
        \begin{tikzpicture}[scale=0.4]
          \begin{scope}
            \lgrid{0}{6}{-1}{5}
            \draw[very thick, forest] (1,3)--(3,1)--(5,3);
            \node[charcoal, font=\scriptsize] at (3,-2.4){\textbf{I}};
          \end{scope}
          \begin{scope}[xshift=9cm]
            \lgrid{-5}{1}{-1}{4}
            \draw[very thick, forest] (-4,0)--(-2,2)--(0,0);
            \node[charcoal, font=\scriptsize] at (-2,-2.4){\textbf{II}};
          \end{scope}
          \begin{scope}[xshift=18cm]
            \lgrid{-3}{3}{-1}{5}
            \draw[very thick, forest] (-2,4)--(0,0)--(2,4);
            \node[charcoal, font=\scriptsize] at (0,-2.4){\textbf{III}};
          \end{scope}
        \end{tikzpicture}
        \end{center}
        $y=2|x|$ is \ans{III}. \quad $y=|x-3|+1$ is \ans{I}. \quad $y=-|x+2|+2$ is \ans{II}.
  \item \textbf{Describe the transformation} from $y=|x|$ to $y=-|x-1|+3$, then give the new vertex.
        \ansline{Shift right $1$, reflect over the $x$-axis (open down), and shift up $3$.}
        Vertex: $({\color{keyred}\mathbf{1}},\,{\color{keyred}\mathbf{3}})$.
  \item \textbf{Evaluate.} For $g(x)=3|x-2|-1$: \quad $g(2)={\color{keyred}\mathbf{-1}}$ \qquad $g(4)={\color{keyred}\mathbf{5}}$.
\end{enumerate}
\end{tcolorbox}

\begin{tcolorbox}[colback=white, colframe=black!40, title=\textbf{Tier E --- Extension},
    fonttitle=\bfseries, arc=1.5mm, left=3mm, right=3mm, top=2mm, bottom=2mm, breakable]
\begin{enumerate}[leftmargin=1.7em, itemsep=10pt]
  \item \textbf{Model it.} A cyclist rides straight toward a water station, reaches it $4$ minutes in,
        and continues past. Her distance from the station (in meters) is $d(t)=80\,|t-4|$.
        \begin{enumerate}[label=(\alph*), leftmargin=1.8em, itemsep=5pt]
          \item What is the vertex of this graph, and what does it mean in context?
                \ansline{Vertex $(4,0)$: at $t=4$ min she is \emph{at} the station, distance $0$ (closest).}
          \item How far from the station is she at $t=0$ and at $t=7$?
                \ans{$d(0)=80(4)=320$ m; \; $d(7)=80(3)=240$ m.}
        \end{enumerate}
  \item \textbf{Build the equation from a graph.} The V below opens down.
        \begin{center}
        \begin{tikzpicture}[scale=0.5]
          \lgrid{-4}{2}{-1}{5}
          \draw[very thick, forest] (-3,0)--(-1,4)--(1,0);
          \fill[charcoal] (-1,4) circle (3.5pt);
        \end{tikzpicture}
        \end{center}
        Vertex $({\color{keyred}\mathbf{-1}},\,{\color{keyred}\mathbf{4}})$; using a second point,
        $a={\color{keyred}\mathbf{-2}}$; equation $y=\ans{$-2|x+1|+4$}$.
  \item \textbf{Reason.} For $y=a|x-h|+k$, explain how you know from $a$ alone whether the vertex is a
        \emph{maximum} or a \emph{minimum}, and state the range in each case.
        \ansline{If $a>0$ the V opens up, so the vertex is a \emph{minimum} and the range is $y\ge k$.}
        \ansline{If $a<0$ the V opens down, so the vertex is a \emph{maximum} and the range is $y\le k$.}
\end{enumerate}
\end{tcolorbox}

\begin{teachernote}
Tier~A item~2 and Tier~E items~2--3 are the conceptual core --- push students to locate the vertex by
setting the inside of the bars to zero, and to read $a$ as ``sign $=$ direction, size $=$ width.''
Common slips: in A1(b), forgetting that $a<0$ makes $0$ a \emph{maximum}; in A3, reading $|x-1|$ as a
left shift; in E2, computing $a$ with the wrong run or dropping the negative (the graph opens down, so
$a<0$). For E1, stress that the vertex $(4,0)$ is the \emph{closest} moment because a distance can't be
negative --- the same distance idea that makes every V.
\end{teachernote}

\end{document}
